% Shared preamble for the TikZ figures; each figure starts with
%   \documentclass[border=4pt]{standalone}
%   % Shared preamble for the TikZ figures; each figure starts with
%   \documentclass[border=4pt]{standalone}
%   % Shared preamble for the TikZ figures; each figure starts with
%   \documentclass[border=4pt]{standalone}
%   % Shared preamble for the TikZ figures; each figure starts with
%   \documentclass[border=4pt]{standalone}
%   \input{_preamble}
% Coordinates are px on a 1000px-wide canvas with y growing downwards (the
% `px` style), and font sizes are px * 0.75. Figures are drawn for a light
% background; style.css inverts them in dark mode.
\def\pgfsysdriver{pgfsys-dvisvgm.def}
\usepackage[T1]{fontenc}
\usepackage{lmodern}
\usepackage[scaled]{helvet}
\renewcommand{\familydefault}{\sfdefault}
\usepackage{tikz}
\usetikzlibrary{arrows.meta, calc}

\definecolor{ink}{HTML}{222222}
\definecolor{muted}{HTML}{6A6A60}
\definecolor{edge}{HTML}{808080}
\definecolor{accent}{HTML}{D64A3C}
\definecolor{accentbg}{HTML}{FBE4E1}
\definecolor{unit}{HTML}{4A6FA5}
\definecolor{latentbg}{HTML}{E8EEF6}
\definecolor{observedbg}{HTML}{C9C9C9}
\definecolor{observed}{HTML}{444444}
\definecolor{panel}{HTML}{F0F0F0}

% Semantic math colors, same names and light values as the KaTeX macros in
% build.js / --m-* in style.css (the dark-mode invert keeps the hue)
\definecolor{mq}{HTML}{1A6ABF}     % approximation q
\definecolor{mpost}{HTML}{B5402A}  % true posterior
\definecolor{mprior}{HTML}{2E7D32} % prior
\definecolor{mlik}{HTML}{B8680A}   % likelihood
\definecolor{mev}{HTML}{6F42C1}    % evidence
\definecolor{mpar}{HTML}{AD1457}   % variational parameters mu, sigma
\newcommand{\cq}[1]{\textcolor{mq}{#1}}
\newcommand{\cpost}[1]{\textcolor{mpost}{#1}}
\newcommand{\cprior}[1]{\textcolor{mprior}{#1}}
\newcommand{\clik}[1]{\textcolor{mlik}{#1}}
\newcommand{\cev}[1]{\textcolor{mev}{#1}}
\newcommand{\cpar}[1]{\textcolor{mpar}{#1}}

\newcommand{\fsBody}{\fontsize{11.25}{13.5}\selectfont}   % 15px
\newcommand{\fsLabel}{\fontsize{10.5}{12.6}\selectfont}   % 14px
\newcommand{\fsSmall}{\fontsize{9.75}{11.7}\selectfont}   % 13px
\newcommand{\fsTiny}{\fontsize{9}{10.8}\selectfont}       % 12px

\tikzset{
  px/.style={x=0.75pt, y=-0.75pt, font=\fsBody, text=ink, line join=round},
  edge/.style={draw=edge, line width=1.125pt},
  % network nodes: observed inputs/outputs and hidden units
  io/.style={circle, draw=edge, fill=white, line width=1.125pt, minimum size=33pt,
             inner sep=0pt, text=muted},
  hidden/.style={circle, fill=unit, minimum size=27pt, inner sep=0pt},
  caption/.style={font=\fsBody\bfseries},
}

% \gauss[color]{(x,y)}{mu}{sigma}{half width}{peak}: N(mu, sigma^2) over [-3, 3]
% stretched to +-half width (px), above a baseline at (x,y) with a tick at 0.
% The peak is `peak`/sigma px tall, capped at 26px. A white backing hides
% the edge the glyph sits on.
\newcommand{\gauss}[6][accent]{%
  \begin{scope}[shift={#2}]
    \pgfmathsetmacro{\gaussh}{min(26, #6/#4)}
    \fill[white] (-#5-4, 6) rectangle (#5+4, -\gaussh-2);
    \draw[edge, line width=0.75pt] (-#5, 0) -- (#5, 0) (0, 0) -- (0, 4);
    \draw[#1, line width=1.35pt, domain=-3:3, samples=241]
      plot ({\x*#5/3}, {-\gaussh*exp(-(\x-(#3))^2/(2*(#4)^2))});
  \end{scope}}

% Coordinates are px on a 1000px-wide canvas with y growing downwards (the
% `px` style), and font sizes are px * 0.75. Figures are drawn for a light
% background; style.css inverts them in dark mode.
\def\pgfsysdriver{pgfsys-dvisvgm.def}
\usepackage[T1]{fontenc}
\usepackage{lmodern}
\usepackage[scaled]{helvet}
\renewcommand{\familydefault}{\sfdefault}
\usepackage{tikz}
\usetikzlibrary{arrows.meta, calc}

\definecolor{ink}{HTML}{222222}
\definecolor{muted}{HTML}{6A6A60}
\definecolor{edge}{HTML}{808080}
\definecolor{accent}{HTML}{D64A3C}
\definecolor{accentbg}{HTML}{FBE4E1}
\definecolor{unit}{HTML}{4A6FA5}
\definecolor{latentbg}{HTML}{E8EEF6}
\definecolor{observedbg}{HTML}{C9C9C9}
\definecolor{observed}{HTML}{444444}
\definecolor{panel}{HTML}{F0F0F0}

% Semantic math colors, same names and light values as the KaTeX macros in
% build.js / --m-* in style.css (the dark-mode invert keeps the hue)
\definecolor{mq}{HTML}{1A6ABF}     % approximation q
\definecolor{mpost}{HTML}{B5402A}  % true posterior
\definecolor{mprior}{HTML}{2E7D32} % prior
\definecolor{mlik}{HTML}{B8680A}   % likelihood
\definecolor{mev}{HTML}{6F42C1}    % evidence
\definecolor{mpar}{HTML}{AD1457}   % variational parameters mu, sigma
\newcommand{\cq}[1]{\textcolor{mq}{#1}}
\newcommand{\cpost}[1]{\textcolor{mpost}{#1}}
\newcommand{\cprior}[1]{\textcolor{mprior}{#1}}
\newcommand{\clik}[1]{\textcolor{mlik}{#1}}
\newcommand{\cev}[1]{\textcolor{mev}{#1}}
\newcommand{\cpar}[1]{\textcolor{mpar}{#1}}

\newcommand{\fsBody}{\fontsize{11.25}{13.5}\selectfont}   % 15px
\newcommand{\fsLabel}{\fontsize{10.5}{12.6}\selectfont}   % 14px
\newcommand{\fsSmall}{\fontsize{9.75}{11.7}\selectfont}   % 13px
\newcommand{\fsTiny}{\fontsize{9}{10.8}\selectfont}       % 12px

\tikzset{
  px/.style={x=0.75pt, y=-0.75pt, font=\fsBody, text=ink, line join=round},
  edge/.style={draw=edge, line width=1.125pt},
  % network nodes: observed inputs/outputs and hidden units
  io/.style={circle, draw=edge, fill=white, line width=1.125pt, minimum size=33pt,
             inner sep=0pt, text=muted},
  hidden/.style={circle, fill=unit, minimum size=27pt, inner sep=0pt},
  caption/.style={font=\fsBody\bfseries},
}

% \gauss[color]{(x,y)}{mu}{sigma}{half width}{peak}: N(mu, sigma^2) over [-3, 3]
% stretched to +-half width (px), above a baseline at (x,y) with a tick at 0.
% The peak is `peak`/sigma px tall, capped at 26px. A white backing hides
% the edge the glyph sits on.
\newcommand{\gauss}[6][accent]{%
  \begin{scope}[shift={#2}]
    \pgfmathsetmacro{\gaussh}{min(26, #6/#4)}
    \fill[white] (-#5-4, 6) rectangle (#5+4, -\gaussh-2);
    \draw[edge, line width=0.75pt] (-#5, 0) -- (#5, 0) (0, 0) -- (0, 4);
    \draw[#1, line width=1.35pt, domain=-3:3, samples=241]
      plot ({\x*#5/3}, {-\gaussh*exp(-(\x-(#3))^2/(2*(#4)^2))});
  \end{scope}}

% Coordinates are px on a 1000px-wide canvas with y growing downwards (the
% `px` style), and font sizes are px * 0.75. Figures are drawn for a light
% background; style.css inverts them in dark mode.
\def\pgfsysdriver{pgfsys-dvisvgm.def}
\usepackage[T1]{fontenc}
\usepackage{lmodern}
\usepackage[scaled]{helvet}
\renewcommand{\familydefault}{\sfdefault}
\usepackage{tikz}
\usetikzlibrary{arrows.meta, calc}

\definecolor{ink}{HTML}{222222}
\definecolor{muted}{HTML}{6A6A60}
\definecolor{edge}{HTML}{808080}
\definecolor{accent}{HTML}{D64A3C}
\definecolor{accentbg}{HTML}{FBE4E1}
\definecolor{unit}{HTML}{4A6FA5}
\definecolor{latentbg}{HTML}{E8EEF6}
\definecolor{observedbg}{HTML}{C9C9C9}
\definecolor{observed}{HTML}{444444}
\definecolor{panel}{HTML}{F0F0F0}

% Semantic math colors, same names and light values as the KaTeX macros in
% build.js / --m-* in style.css (the dark-mode invert keeps the hue)
\definecolor{mq}{HTML}{1A6ABF}     % approximation q
\definecolor{mpost}{HTML}{B5402A}  % true posterior
\definecolor{mprior}{HTML}{2E7D32} % prior
\definecolor{mlik}{HTML}{B8680A}   % likelihood
\definecolor{mev}{HTML}{6F42C1}    % evidence
\definecolor{mpar}{HTML}{AD1457}   % variational parameters mu, sigma
\newcommand{\cq}[1]{\textcolor{mq}{#1}}
\newcommand{\cpost}[1]{\textcolor{mpost}{#1}}
\newcommand{\cprior}[1]{\textcolor{mprior}{#1}}
\newcommand{\clik}[1]{\textcolor{mlik}{#1}}
\newcommand{\cev}[1]{\textcolor{mev}{#1}}
\newcommand{\cpar}[1]{\textcolor{mpar}{#1}}

\newcommand{\fsBody}{\fontsize{11.25}{13.5}\selectfont}   % 15px
\newcommand{\fsLabel}{\fontsize{10.5}{12.6}\selectfont}   % 14px
\newcommand{\fsSmall}{\fontsize{9.75}{11.7}\selectfont}   % 13px
\newcommand{\fsTiny}{\fontsize{9}{10.8}\selectfont}       % 12px

\tikzset{
  px/.style={x=0.75pt, y=-0.75pt, font=\fsBody, text=ink, line join=round},
  edge/.style={draw=edge, line width=1.125pt},
  % network nodes: observed inputs/outputs and hidden units
  io/.style={circle, draw=edge, fill=white, line width=1.125pt, minimum size=33pt,
             inner sep=0pt, text=muted},
  hidden/.style={circle, fill=unit, minimum size=27pt, inner sep=0pt},
  caption/.style={font=\fsBody\bfseries},
}

% \gauss[color]{(x,y)}{mu}{sigma}{half width}{peak}: N(mu, sigma^2) over [-3, 3]
% stretched to +-half width (px), above a baseline at (x,y) with a tick at 0.
% The peak is `peak`/sigma px tall, capped at 26px. A white backing hides
% the edge the glyph sits on.
\newcommand{\gauss}[6][accent]{%
  \begin{scope}[shift={#2}]
    \pgfmathsetmacro{\gaussh}{min(26, #6/#4)}
    \fill[white] (-#5-4, 6) rectangle (#5+4, -\gaussh-2);
    \draw[edge, line width=0.75pt] (-#5, 0) -- (#5, 0) (0, 0) -- (0, 4);
    \draw[#1, line width=1.35pt, domain=-3:3, samples=241]
      plot ({\x*#5/3}, {-\gaussh*exp(-(\x-(#3))^2/(2*(#4)^2))});
  \end{scope}}

% Coordinates are px on a 1000px-wide canvas with y growing downwards (the
% `px` style), and font sizes are px * 0.75. Figures are drawn for a light
% background; style.css inverts them in dark mode.
\def\pgfsysdriver{pgfsys-dvisvgm.def}
\usepackage[T1]{fontenc}
\usepackage{lmodern}
\usepackage[scaled]{helvet}
\renewcommand{\familydefault}{\sfdefault}
\usepackage{tikz}
\usetikzlibrary{arrows.meta, calc}

\definecolor{ink}{HTML}{222222}
\definecolor{muted}{HTML}{6A6A60}
\definecolor{edge}{HTML}{808080}
\definecolor{accent}{HTML}{D64A3C}
\definecolor{accentbg}{HTML}{FBE4E1}
\definecolor{unit}{HTML}{4A6FA5}
\definecolor{latentbg}{HTML}{E8EEF6}
\definecolor{observedbg}{HTML}{C9C9C9}
\definecolor{observed}{HTML}{444444}
\definecolor{panel}{HTML}{F0F0F0}

% Semantic math colors, same names and light values as the KaTeX macros in
% build.js / --m-* in style.css (the dark-mode invert keeps the hue)
\definecolor{mq}{HTML}{1A6ABF}     % approximation q
\definecolor{mpost}{HTML}{B5402A}  % true posterior
\definecolor{mprior}{HTML}{2E7D32} % prior
\definecolor{mlik}{HTML}{B8680A}   % likelihood
\definecolor{mev}{HTML}{6F42C1}    % evidence
\definecolor{mpar}{HTML}{AD1457}   % variational parameters mu, sigma
\newcommand{\cq}[1]{\textcolor{mq}{#1}}
\newcommand{\cpost}[1]{\textcolor{mpost}{#1}}
\newcommand{\cprior}[1]{\textcolor{mprior}{#1}}
\newcommand{\clik}[1]{\textcolor{mlik}{#1}}
\newcommand{\cev}[1]{\textcolor{mev}{#1}}
\newcommand{\cpar}[1]{\textcolor{mpar}{#1}}

\newcommand{\fsBody}{\fontsize{11.25}{13.5}\selectfont}   % 15px
\newcommand{\fsLabel}{\fontsize{10.5}{12.6}\selectfont}   % 14px
\newcommand{\fsSmall}{\fontsize{9.75}{11.7}\selectfont}   % 13px
\newcommand{\fsTiny}{\fontsize{9}{10.8}\selectfont}       % 12px

\tikzset{
  px/.style={x=0.75pt, y=-0.75pt, font=\fsBody, text=ink, line join=round},
  edge/.style={draw=edge, line width=1.125pt},
  % network nodes: observed inputs/outputs and hidden units
  io/.style={circle, draw=edge, fill=white, line width=1.125pt, minimum size=33pt,
             inner sep=0pt, text=muted},
  hidden/.style={circle, fill=unit, minimum size=27pt, inner sep=0pt},
  caption/.style={font=\fsBody\bfseries},
}

% \gauss[color]{(x,y)}{mu}{sigma}{half width}{peak}: N(mu, sigma^2) over [-3, 3]
% stretched to +-half width (px), above a baseline at (x,y) with a tick at 0.
% The peak is `peak`/sigma px tall, capped at 26px. A white backing hides
% the edge the glyph sits on.
\newcommand{\gauss}[6][accent]{%
  \begin{scope}[shift={#2}]
    \pgfmathsetmacro{\gaussh}{min(26, #6/#4)}
    \fill[white] (-#5-4, 6) rectangle (#5+4, -\gaussh-2);
    \draw[edge, line width=0.75pt] (-#5, 0) -- (#5, 0) (0, 0) -- (0, 4);
    \draw[#1, line width=1.35pt, domain=-3:3, samples=241]
      plot ({\x*#5/3}, {-\gaussh*exp(-(\x-(#3))^2/(2*(#4)^2))});
  \end{scope}}
